\documentclass[a4paper,11pt,fleqn]{article}%fleqn pour aligner les equations à gauche
\usepackage{../../Styles/Cours}

\newcommand{\chnum}{Chapitre 2}
\newcommand{\chtitle}{Bouillie Bordelaise}
\newcommand{\dtype}{TP}
  
\begin{document}
\maketitle

\noindent
\fbox{\begin{minipage}{\linewidth}
\textbf{Document 1 : La bouillie bordelaise}\\
La bouillie bordelaise est un produit de traitement fongicide efficace et autorisé en agriculture biologique. Elle est largement utilisée pour le traitement des plantes, légumes ou fruitiers du jardin. Le respect des doses et de l’utilisation de ce produit est néanmoins nécessaire pour ne pas contaminer la nature.\\
\textbf{Utilisation} : fongicide, algicide\\
\textbf{Composition} : eau, sulfate de cuivre et chaux\\
\textbf{Principales plantes concernées} : vigne, fruitiers, pomme de terre, tomates\\
\end{minipage}}
\vspace{.3cm}

\noindent
\fbox{\begin{minipage}{\linewidth}
\textbf{Document 2 : Utilisation de la bouillie bordelaise}\\
Il s’agit d’un fongicide de couleur bleue à base de sulfate de cuivre et de chaux. Elle est utilisée de préférence en pulvérisation et sert à lutter contre la plupart des formes de maladies cryptogamiques, c’est à dire les champignons et couramment utilisée contre le mildiou, la tavelure, la cloque des fruits, les chancres ou encore la moniliose.\\
\end{minipage}}
\vspace{.3cm}

\noindent
\fbox{\begin{minipage}{\linewidth}
\textbf{Document 3 : Préparation des solutions filles}\\
Pour évaluer la concentration d'une solution de sulfate de cuivre (\ce{Cu^{2+}} ; \ce{SO4^{2-}}) de concentration en masse notée $\gamma$ , il faut disposer d’une gamme étalon.\\
On dispose d’une solution mère $S_0$ de concentration connue égale à \num{250} \unit{\gram\per\liter} à partir de laquelle sont préparées 5 solutions filles notées $S_1$, $S_2$, $S_3$, $S_4$ et $S_5$.\\
On dispose de fioles, de pipettes jaugées et graduées de différents volumes. Chaque groupe expliquera la préparation d’une solution fille. Il faudra déterminer la fiole et la pipette à utiliser pour préparer chaque solution.\\
\begin{tabular}{|>{\centering}m{3cm}|>{\centering}m{2.5cm}|>{\centering}m{2.5cm}|>{\centering}m{2.5cm}|>{\centering}m{2.5cm}|>{\centering\arraybackslash}m{2.5cm}|}
\hline Solution & $S_1$ & $S_2$ & $S_3$ & $S_4$ & $S_5$\\
\hline Concentration (\unit{\gram\per\liter})& 25& 50& 75& 100& 125\\
\hline $V_0$ à prélever (mL)& & & & & \rule[1cm]{0cm}{0cm}\\
\hline $V_{solution}$ (\unit{\milli\liter})&50,0 &50,0 &50,0 &50,0 &50,0\\
\hline
\end{tabular}
\end{minipage}}
\vspace{.3cm}

\noindent
\fbox{\begin{minipage}{\linewidth}
\textbf{Document 4 : Méthode de préparation par dilution}\\
- Prélever un volume $V_0$ de la solution mère $S_0$ à l’aide d’une pipette jaugée ou graduée et l’introduire dans une fiole jaugée de volume $V$.\\
- Compléter la fiole jaugée avec de l’eau distillée jusqu’au trait de jauge puis homogénéiser la solution.
\end{minipage}}\\

\textbf{Questions :}
\begin{enumerate}
	\item Ecrire un protocole de préparation pour chaque solution fille en précisant le matériel utilisé.
	\item Compléter le tableau pour toutes les solutions en détaillant les calculs dans chaque cas.
	\item Comparer la solution commerciale aux solutions de l’échelle de teinte et en déduire un encadrement
de sa concentration.
	\item L’étiquette de la bouillie bordelaise commerciale mentionne « pour 100 mL de solution, permanganate de potassium : 6,25 g ». Calculer la concentration en masse de sulfate de cuivre de la solution de bouillie bordelaise.
	\item La valeur calculée est elle en accord avec celle de la concentration estimée à l’aide de l’échelle
de teinte ?
\end{enumerate}
\end{document}